% s07 Reflection positivity I: the modified box
\section{Reflection positivity I: the modified box}

King's Neumann--Dirichlet cube is not reflection symmetric about any plane, and the natural symmetrisation of its photon form is not reflection positive on gauge-invariant functionals (\cref{rem:s07:R1}). Reflection positivity is therefore proved first for the modified box state of \cite[(3.38)]{K2} and then carried to the cube states by a two-sided comparison on the cone: the cube $C_J$ sits inside the box (\cref{lem:s07:L}), and the box sits inside a larger cube once a margin proportional to $\lmu^{-1}$ is allowed (\cref{lem:s07:U}). The passage to the limits is in \cref{thm:s08:OS2C}. Throughout, $\eps=\eps_k$, the lattice is $\Lk\subset\R^d$ with $d=3$ (the certificates of \cref{lem:s07:R} are also given for $d=2$), a \emph{lattice box} is a product $\Omega=\prod_\mu\{a_\mu,a_\mu+\eps,\dots,b_\mu\}\subset\Lk$, $\Omega^*$ is the set of its internal bonds (both ends in $\Omega$) and $\Omega^{**}$ the set of its plaquettes (all four corners in $\Omega$).

\subsection{The lattice reflection and the modified box}

\begin{definition}[mid-planes and the lattice reflection]\label{def:s07:reflection}
The \emph{mid-planes} are the planes $\{x_\nu=c\}$ with $c\in D\defeq\bigcup_{k\ge0}\eps_k(\Z+\frac12)$. If $c\in\eps_{k_0}(\Z+\frac12)$ then $c\in\eps_k(\Z+\frac12)$ for every $k\ge k_0$, because $5^j(m+\frac12)=5^jm+\frac{5^j-1}2+\frac12$ and $5^j$ is odd; the same set $D$ is obtained from any tail $k\ge k_1$ of the spacings. $D$ is dense in $\R$ but not all of $\R$.

Fix $\nu$, a mid-plane $c\in\eps_k(\Z+\frac12)$ and a lattice box $\Omega\subset\Lk$ with $\rho\Omega=\Omega$, $\rho=\rho_{\nu,c}$ the reflection of \cref{def:s02:sources}. On the fields put $(\hat\rho\phi)(x)\defeq\phi(\rho x)$ and $(\hat\rho A)(b)\defeq A(\rho b)$, the image bond taken with its orientation: $(\hat\rho A)_\alpha(x)=A_\alpha(\rho x)$ for $\alpha\ne\nu$ and $(\hat\rho A)_\nu(x)=-A_\nu(\rho x-\eps e_\nu)$. Then $\hat\rho$ is a signed permutation of the variables attached to $\Omega$ and $\Omega^*$ (it preserves Lebesgue measure and maps $\Omega^*$, $\Omega^{**}$ to themselves); the \emph{crossing bonds} $B_0\defeq\{\langle x^-,x^-+\eps e_\nu\rangle:\ x^-_\nu=c-\eps/2\}$ are mapped to their reversals, $(\hat\rho A)(b)=-A(b)$ for $b\in B_0$. Let $\Omega_\pm\defeq\{x\in\Omega:\ \pm(x_\nu-c)>0\}$ (no site lies on the plane), $B_+$ the bonds with both ends in $\Omega_+$, $\mathfrak B_+$ the Borel functions of $(\phi|_{\Omega_+},A|_{B_+})$, and
\[
\theta G\defeq\overline{G\circ\hat\rho}\qquad(G\in\mathfrak B_+).
\]
\end{definition}

On the smeared fields, $F_{\Lk}(f)\circ\hat\rho=F_{\Lk}(\rho^*f)$ and $W^a_{\Lk}(h)\circ\hat\rho=W^a_{\Lk}(h\circ\rho)$. Indeed, $F(p)\circ\hat\rho=\sigma_pF(\rho p)$ with $\sigma_p=-1$ exactly when $\nu\in\{\alpha,\beta\}$ (the orientation of the image plaquette, whose base point is shifted back one layer), the centres satisfy $c(\rho p)=\rho c(p)$, and $(\rho^*f)_{c(\rho p)}(e_\alpha,e_\beta)=f_{c(p)}(\rho e_\alpha,\rho e_\beta)=\sigma_pf_{c(p)}(e_\alpha,e_\beta)$; reindexing $p\mapsto\rho p$ gives the first identity. The scalar case is a relabelling of sites, and the $c$-numbers $2\sW_\eps\sS(h)$ and $a_k\sS(h)$ are invariant because $\rho$ permutes $\Lk$ ($2(x_\nu-c)\in\eps\Z$). Consequently, for families $\zeta_i$ whose lattice observables $\Theta(\zeta_i)$ lie in $\mathfrak B_+$,
\begin{equation}\label{eq:s07:thetaTheta}
\theta(\Theta(\zeta_i))\,\Theta(\zeta_j)=\Theta(\theta\zeta_i^*\sqcup\zeta_j),\qquad m(\theta\zeta_i^*\sqcup\zeta_j)=\overline{m(\zeta_i)}\,m(\zeta_j)
\end{equation}
for the scheme multiplier $m(\zeta)\defeq e^{-\sum_jw_ja_k\sS(h_j)}$ ($a_k$ real), so that for any state $\av{\cdot}$ on the box
\begin{equation}\label{eq:s07:matrix}
\sum_{i,j}\bar\gamma_i\gamma_j\,m(\theta\zeta_i^*\sqcup\zeta_j)\av{\Theta(\theta\zeta_i^*\sqcup\zeta_j)}=\sum_{i,j}\overline{(\gamma_im(\zeta_i))}\,(\gamma_jm(\zeta_j))\,\big\langle\theta(\Theta(\zeta_i))\,\Theta(\zeta_j)\big\rangle .
\end{equation}
\emph{Support.} If the sources of $\zeta_i$ are supported in $\{x_\nu\ge c+\delta\}$ with $\delta>0$, then for $\eps_k<2\delta$ every site carrying a sampled $h$ and every corner of every plaquette carrying a sampled $f$ (centre in the support, corners at $x_\nu\ge c+\delta-\eps_k/2>c$) lies in $\Omega_+$, so $\Theta(\zeta_i)\in\mathfrak B_+$. The splitting $h=h^+-h^-$ of \cref{def:s08:enlarged} is taken with cutoffs supported in $\{x_\nu\ge c+\delta/2\}$, so that $h^\pm$ are supported in $\{x_\nu>c\}$ as well.

\begin{definition}[the modified box state]\label{def:s07:box}
Let $\Omega\subset\Lk$ be a lattice box, $F(p)=\eps^{-1}\sum_{b\in\partial p}A(b)$, $\Omega(y)$ the indicator of $\Omega$, and $x^*$ the nearest neighbours of $x$ in $\Lk$. King's modified action \cite[(3.38)]{K2} on $\Omega$ is
\begin{align*}
S^\eps_\Omega(A,\phi)&=\tfrac12\sum_{p\in\Omega^{**}}\eps^d|F(p)|^2+\tfrac12\lmu^2\sum_{b\in\Omega^*}\eps^d|A(b)|^2+\tfrac12\sum_{x\in\Omega}\eps^d\Big|\sum_{y\in x^*}\eps^{-1}\Omega(y)A(\langle x,y\rangle)\Big|^2\\
&\quad+\tfrac12\sum_{\langle x,y\rangle}\eps^{d-2}\big[(\Omega r)(x)^2+(\Omega r)(y)^2-2(\Omega r)(x)(\Omega r)(y)\cos(\vartheta_x-e\eps A(\langle x,y\rangle)-\vartheta_y)\big]\\
&\quad+\sum_{x\in\Omega}\eps^d\big[\tfrac12m^2r_x^2+\lambda r_x^4+\tfrac12\dm_\eps\,r_x^2\big],
\end{align*}
with the \emph{constant} free-boundary counterterm $\dm_\eps$ of \cref{def:s02:Kcube} at every site, and with the free Wick constant $\sW_\eps$ in the Wick observables. (King prints a site-dependent counterterm in the last term; it is not used here, and the statements below are made for the constant one.) The photon variables are the internal bonds $\Omega^*$ only. The matter part is exactly the Dirichlet matter of the Neumann--Dirichlet state on $X=\Omega$: cosines on the internal bonds, and the diagonal $\frac12\eps^{d-2}r_x^2$ for every bond leaving $\Omega$. The \emph{modified box state} $\langle\cdot\rangle^{(3.38)}_{k,\Omega}$ is the normalised measure $e^{-S^\eps_\Omega}\prod_{x\in\Omega}d^2\phi_x\prod_{b\in\Omega^*}dA_b$; it is a BFS measure (the Gaussian is non-degenerate since $\lmu>0$, the single-site measures have all exponential moments since $\lambda>0$). We write $Z^{a,(3.38)}_{k,\Omega}(\zeta)\defeq m(\zeta)\langle\Theta(\zeta)\rangle^{(3.38)}_{k,\Omega}$ with the same observables and multiplier as in \cref{def:s02:sources}.

The third term is the graph divergence of $(\Omega,\Omega^*)$: with $d_0:\R^\Omega\to\R^{\Omega^*}$, $(d_0\chi)(\langle x,x+\eps e_\alpha\rangle)=\eps^{-1}(\chi(x+\eps e_\alpha)-\chi(x))$, one has $\sum_{y\in x^*}\Omega(y)A(\langle x,y\rangle)=-\eps(d_0^TA)(x)$ (only internal bonds contribute). With $d_1A\defeq(\eps F(p))_{p\in\Omega^{**}}$, the \emph{unitary form} on the box and the modified form are
\[
Q^u_\Omega(A)\defeq\tfrac12\sum_{p\in\Omega^{**}}\eps^d|F(p)|^2+\tfrac12\lmu^2\sum_{b\in\Omega^*}\eps^d|A(b)|^2,\qquad Q^{(3.38)}_\Omega=Q^u_\Omega+\tfrac12\eps^d\|d_0^TA\|^2 ,
\]
\ie the modified form is the unitary form plus a gauge-fixing term depending on $A$ only through $d_0^TA$. The modified action is invariant under $\hat\rho$ for every mid-plane $c$ with $\rho\Omega=\Omega$: every term is written through $A(\langle x,y\rangle)$ and nearest-neighbour pairs, which $\rho$ permutes, $(d_0^T\hat\rho A)(x)=(d_0^TA)(\rho x)$, and $\dm_\eps$ is site independent.
\end{definition}

\emph{The boxes.} Fix $c\in D$, $k\ge k_0(c)$, $J\ge J_0$ such that every source of the enlarged family of \cref{def:s08:enlarged} and its reflection is supported at $\ell^\infty$-distance $>1$ from $\R^3\setminus C_{J_0}$, an integer $n\ge1$, and put $a_k\defeq c+\eps_k/2\in\eps_k\Z$ and
\begin{equation}\label{eq:s07:margin}
\Omega_{J,k}\defeq(C_{J+n}+a_ke_\nu)\cap\Lk,\qquad Mn\ \ge\ |c|+2+R(\lmu),\qquad R(\lmu)\defeq\frac{\pi\sqrt3}{2\lmu}.
\end{equation}
$C_{J+n}\cap\Lk$ is symmetric in direction $\nu$ about $-\eps_k/2$ (sites $x_\nu\in\{-M(J+n),\dots,M(J+n)-\eps_k\}$, and $M(J+n)\in\eps_k\Z$), so $\Omega_{J,k}$ is a lattice box symmetric about $\{x_\nu=c\}$ with an even number $2M(J+n)5^k$ of layers. Under the margin condition of \eqref{eq:s07:margin}, for every $k$: (i) $C_J$ plus one layer is contained in $\Omega_{J,k}$ (this needs $Mn\ge|c|+\frac32\eps_k$); (ii) $C_{J+2n}$ contains every site at $\ell^\infty$-distance $\le R(\lmu)+\eps_k$ from $\Omega_{J,k}$ (this needs $Mn\ge|c|+R(\lmu)+\frac32\eps_k$). We fix such an $n=n(\lmu,c)$, independent of $k$ and $J$. The two margins are exactly the hypotheses of \cref{lem:s07:L,lem:s07:U} below; by \cref{lem:s04:D} they serve King's convention as they stand.

\subsection{Reflection positivity of the modified box}

\begin{lemma}[lattice reflection positivity of the modified box]\label{lem:s07:RPbox}
Let $\lmu>0$, $c$ a mid-plane and $\Omega$ a lattice box with $\rho\Omega=\Omega$, $\rho=\rho_{\nu,c}$.
\begin{enumerate}[label=\textup{(\alph*)},nosep]
\item For every gauge-invariant $G\in\mathfrak B_+$ that is bounded, or of the form $\Theta(\zeta)$ with sources supported in $\Omega_+$, or more generally with $\langle|G|^2\rangle^{(3.38)}_{k,\Omega}<\infty$: $\langle\theta(G)G\rangle^{(3.38)}_{k,\Omega}\ge0$.
\item $(G,G')\mapsto\langle\theta(G')G\rangle^{(3.38)}_{k,\Omega}$ is a Hermitian positive semidefinite form on such $G$; hence $\sum_{i,j}\bar\gamma_i\gamma_j\langle\theta(G_i)G_j\rangle^{(3.38)}_{k,\Omega}\ge0$, and by \eqref{eq:s07:matrix} and the support clause of \cref{def:s07:reflection}, $\sum_{i,j}\bar\gamma_i\gamma_jZ^{a,(3.38)}_{k,\Omega}(\theta\zeta_i^*\sqcup\zeta_j)\ge0$ for families $\zeta_i$ with sources supported in $\{x_\nu\ge c+\delta\}$, $\eps_k<2\delta$.
\end{enumerate}
\end{lemma}

\begin{proof}
\emph{Step 1: gauge-invariant expectations of the modified form equal those of the unitary form.} The graph $(\Omega,\Omega^*)$ is connected, $\ker d_0$ consists of the constants, $\R^{\Omega^*}=\operatorname{ran}d_0\oplus\ker d_0^T$ orthogonally, and $d_1d_0=0$. Write $A=A^T+d_0\chi$ with $d_0^TA^T=0$ and $\chi\perp1$; $(A^T,\chi)\mapsto A$ is a linear bijection with constant Jacobian. Then $Q^u_\Omega(A)=Q^u_\Omega(A^T)+\frac{\lmu^2\eps^d}2\|d_0\chi\|^2$ (orthogonality and $d_1d_0=0$) and $Q^{(3.38)}_\Omega(A)=Q^u_\Omega(A^T)+q(\chi)$ with $q(\chi)=\frac{\lmu^2\eps^d}2\|d_0\chi\|^2+\frac{\eps^d}2\|d_0^Td_0\chi\|^2$; both $\chi$-forms are positive definite on $1^\perp$ because $\lmu>0$. A gauge transformation is $A\mapsto A+d_0\chi$ together with $\vartheta\mapsto\vartheta-e\chi$: the hopping phase $\vartheta_x-e\eps A_b-e(\chi(y)-\chi(x))-\vartheta_y$ equals $(\vartheta_x+e\chi_x)-e\eps A_b-(\vartheta_y+e\chi_y)$ and the Dirichlet diagonals carry no $A$, so for gauge-invariant $G$ the matter integral $\Psi(A)\defeq\int\prod_xd^2\phi_x\,(Ge^{-S_m})(A,\phi)$ satisfies $\Psi(A^T+d_0\chi)=\Psi(A^T)$ (rotate $\phi$ site by site; $\prod d^2\phi$ is invariant). Hence $\int dA\,e^{-Q^{(3.38)}_\Omega}\Psi=\int dA^Te^{-Q^u_\Omega(A^T)}\Psi(A^T)\int d\chi\,e^{-q(\chi)}$, and the same with $Q^u_\Omega$ in place of $Q^{(3.38)}_\Omega$ and a different finite $\chi$-factor; the factors cancel in the normalised expectation. Fubini applies, since $|G|\le Ce^{a\sum_b|A(b)|}e^{C\sum_xr_x^2}$ for $G=\Theta(\zeta)$ (trivially for bounded $G$) is dominated by the Gaussian and by $e^{-\lambda\eps^d r^4}$. This is the mechanism of \cite[Thm.~5.1]{BFS1} and of \cite[Lemma~7.7]{SymPhase}, carried to the box graph.

\emph{Step 2: the unitary box measure is reflection positive on all of $\mathfrak B_+$.} Let $S^u_\Omega\defeq Q^u_\Omega+S_m$, write $\omega=(\omega_+,\omega_-,a)$ with $a=(a_b)_{b\in B_0}$ the crossing-bond variables, and put $\mathsf M\defeq\lmu^2\eps^d$. The terms of $S^u_\Omega$ are of four kinds. Single-site terms (local matter and Dirichlet diagonals), hopping cosines on bonds of $B_\pm$, masses of bonds of $B_\pm$, and plaquettes with all corners in $\Omega_+$ or all in $\Omega_-$ involve one side only; by $\hat\rho$-invariance the minus-side terms are the $\hat\rho$-transport of the plus-side terms. The crossing masses are $\frac{\mathsf M}2a_b^2$, $b\in B_0$. A crossing plaquette $p=(x^-;\alpha,\nu)$, $\alpha\ne\nu$, is present iff $x^-,x^-+\eps e_\alpha\in\Omega$ (then all four corners are in $\Omega$ by symmetry); with $x^+\defeq x^-+\eps e_\nu$, $b=\langle x^-,x^+\rangle$ and $b'=\langle x^-+\eps e_\alpha,x^++\eps e_\alpha\rangle$ it reads $\eps F(p)=\pm[A_\alpha(x^-)+a_{b'}-A_\alpha(x^+)-a_b]$, the overall sign immaterial after squaring. A crossing hopping term is $-\eps^{d-2}r_{x^-}r_{x^+}\cos(\vartheta_{x^-}-e\eps a_b-\vartheta_{x^+})$, $b\in B_0$.

Split each crossing variable: $\int_\R g(a_b)e^{-\frac{\mathsf M}2a_b^2}da_b=\sqrt{2\pi/\mathsf M}\ \E\,g(\alpha_b-\alpha'_b)$ with $\alpha_b,\alpha'_b$ independent $N(0,\frac1{2\mathsf M})$, whose difference is $N(0,\frac1{\mathsf M})$. Put $\xi\defeq(\omega_+,\alpha)$ and $\xi'\defeq(\omega_-^\rho,\alpha')$, where $\omega_-^\rho$ is the $\hat\rho$-image of the minus data transported to the plus side ($(\omega^\rho_-)_\alpha(x^+)=A_\alpha(x^-)$ for the tangential bonds, $\phi^\rho(x^+)=\phi(x^-)$, the $\nu$-bonds of $B_-$ carried with the sign of \cref{def:s07:reflection}); by $\hat\rho$-invariance the minus-side one-sided terms are the plus-side terms evaluated on $\omega_-^\rho$, and the a priori measures of $\xi$ and $\xi'$ coincide. For a crossing plaquette put $X_p(\xi)\defeq A_\alpha(x^+)-\alpha_{b'}+\alpha_b$, a function of plus variables and $\alpha$ (note $A_\alpha(x^+)\in B_+$ since $x^+,x^++\eps e_\alpha\in\Omega_+$), and the same function on $\xi'$, $X_p(\xi')=A_\alpha(x^-)-\alpha'_{b'}+\alpha'_b$; then $\eps F(p)=\pm[X_p(\xi')-X_p(\xi)]$ and
\[
\tfrac{\eps^d}2F(p)^2=\tfrac{\eps^{d-2}}2X_p(\xi)^2+\tfrac{\eps^{d-2}}2X_p(\xi')^2-\eps^{d-2}X_p(\xi)X_p(\xi').
\]
For a crossing hopping term put $\psi_b(\xi)\defeq\phi(x^+)e^{\mathrm ie\eps\alpha_b}$, $\psi_b(\xi')=\phi(x^-)e^{\mathrm ie\eps\alpha'_b}$; then $r_{x^-}r_{x^+}\cos(\vartheta_{x^-}-e\eps(\alpha_b-\alpha'_b)-\vartheta_{x^+})=\Re[\psi_b(\xi')\overline{\psi_b(\xi)}]=\sum_{i=1,2}\psi^i_b(\xi)\psi^i_b(\xi')$. Hence, with $S'$ the action minus the crossing masses,
\[
e^{-S'}=\cG(\xi)\,\cG(\xi')\,\exp\Big(\sum_\iota\gamma_\iota Y_\iota(\xi)Y_\iota(\xi')\Big),\qquad\cG(\xi)\defeq e^{-S_+(\xi)-\frac{\eps^{d-2}}2\sum_pX_p(\xi)^2},
\]
where $S_+$ collects the one-sided plus terms and $\{(\gamma_\iota,Y_\iota)\}=\{(\eps^{d-2},X_p)\}_p\cup\{(\eps^{d-2},\psi^i_b)\}_{b,i}$: all $\gamma_\iota>0$, all $Y_\iota$ real. For $G\in\mathfrak B_+$ we have $G(\omega)=G(\xi)$ and $\theta G(\omega)=\overline{G(\xi')}$, the same function on the transported data ($G$ depends neither on $a$ nor on $\alpha,\alpha'$). Expanding the exponential,
\[
Z\,\langle\theta(G)G\rangle^u_{k,\Omega}=K\sum_{n\in\N^I}\prod_\iota\frac{\gamma_\iota^{n_\iota}}{n_\iota!}\Big|\int d\xi\,G(\xi)\cG(\xi)\prod_\iota Y_\iota(\xi)^{n_\iota}\Big|^2\ \ge0,\qquad K>0,
\]
the two integrals of each term being complex conjugates because measures and functions coincide on the two copies. The interchange of sum and integral is justified by $\sum_n\prod_\iota\gamma_\iota^{n_\iota}|Y_\iota Y'_\iota|^{n_\iota}/n_\iota!\le\exp(\frac12\sum_\iota\gamma_\iota(Y_\iota^2+Y_\iota'^2))$: the plaquette parts cancel against $\cG$ exactly, $|\psi_b|=r_{x^+}$, and $\int|G|e^{-S_+(\xi)+\frac{\eps^{d-2}}2\sum_br_{x^+}^2}d\xi<\infty$ by $\lambda>0$ (a quartic term at every plus site), $\lmu>0$ (a mass on every $B_+$ bond and the $\alpha$-Gaussians) and the growth bound on $G$. This is the mechanism of \cite[Lemma~7.8(ii)]{SymPhase} with one plane in place of a pair and the Dirichlet diagonals as single-site terms; the classical form is the positive cross-term criterion of \cite{GJ,FILS78}. The Hermitian property $\overline{\langle\theta(G')G\rangle}=\langle(G'\circ\hat\rho)\bar G\rangle=\langle G'\,\overline{G\circ\hat\rho}\rangle=\langle\theta(G)G'\rangle$ follows from the $\hat\rho$-invariance of the measure.

\emph{Step 3: combination and $L^2$ extension.} For gauge-invariant $G\in\mathfrak B_+$, bounded or equal to $\Theta(\zeta)$, the product $\theta(G)G$ is gauge invariant ($\hat\rho$ maps gauge transformations to gauge transformations: $\hat\rho\,d_0\chi=d_0(\chi\circ\rho)$), so $\langle\theta(G)G\rangle^{(3.38)}_{k,\Omega}=\langle\theta(G)G\rangle^u_{k,\Omega}\ge0$ by Steps 1--2. For gauge-invariant $G\in\mathfrak B_+$ with $\langle|G|^2\rangle<\infty$ (for instance polynomials in the plus-side smeared fields, \cref{prop:s08:poly}), $G_n\defeq G\,\mathbf 1\{|G|\le n\}$ is bounded, gauge invariant and in $\mathfrak B_+$, so $\langle\theta(G_n)G_n\rangle\ge0$; moreover $|\theta(G_n)G_n|\le|\theta G||G|$ and $\langle|\theta G||G|\rangle\le\langle|G|^2\rangle^{1/2}\langle|\theta G|^2\rangle^{1/2}=\langle|G|^2\rangle$ by Cauchy--Schwarz and $\hat\rho$-invariance, so dominated convergence gives (a). The modified measure is itself $\hat\rho$-invariant (\cref{def:s07:box}), so the computation of the Hermitian property in Step 2 holds verbatim for it on $L^2$; the form is sesquilinear, Hermitian and nonnegative on the diagonal, hence positive semidefinite (apply (a) to $G=\sum_i\gamma_iG_i$). With \eqref{eq:s07:matrix} and the support clause this is (b).
\end{proof}

\begin{remark}[on King's Proposition 3.3]
\Cref{lem:s07:RPbox} is the statement of \cite[Prop.~3.3]{K2}, which King derives from \cite[Cor.~5.4]{BFS1}. That corollary (with \cite[Lemma~5.2, Thm.~5.3]{BFS1}) concerns the lattice actions of \cite[\S5]{BFS1}, whose photon Gaussian is translation invariant with a reflection-positive boundary term; King's finite-box Gaussian with its masked divergence is not one of them, the obstruction being the box geometry of the gauge-fixing term. So his Proposition 3.3 is true, and it is proved here rather than derived from \cite[Cor.~5.4]{BFS1}. The gauge-fixing term $\|d_0^TA\|^2$ is not itself compatible with the mechanism of Step 2 (at a site next to the plane it produces cross terms $Y(\xi)Z(\xi')+Z(\xi)Y(\xi')$ of indefinite sign); on the torus it drops out for gauge-invariant observables, on King's cube it does not reduce to a divergence at the boundary (\cref{lem:s07:B}), and this is why the photon form has to be changed before reflecting (\cref{rem:s07:R1}).
\end{remark}

\subsection{The same box: the boundary identity and the reversed inequality}

In this subsection we use lattice units, $\eps=1$ (every term of the three photon forms carries the same weight $\frac12\eps^{d-2}$), and write $\mu\defeq\lmu\eps$. For a lattice box $\Omega$ and $A\in\R^{\Omega^*}$ let $K^{\rm int}_\Omega(A)\defeq\sum_\mu\sum_\nu\sum|A_\mu(x+e_\nu)-A_\mu(x)|^2$, the sum over pairs with both bonds $(x,\mu),(x+e_\nu,\mu)$ in $\Omega^*$ (the componentwise Neumann gradient form on the internal bonds), and $N_\partial(A)\defeq\sum|A_b|^2$, the sum over the normal crossing pairs: for $b=(x,\mu)\in\Omega^*$ one term for each of $(x\pm e_\mu,\mu)$ that is not in $\Omega^*$. Split $N_\partial=N^{\rm up}_\partial+N^{\rm low}_\partial$ according to whether the missing partner is $(x+e_\mu,\mu)$ (\ie $x+2e_\mu\notin\Omega$) or $(x-e_\mu,\mu)$ (\ie $x-e_\mu\notin\Omega$).

\begin{lemma}[the boundary identity]\label{lem:s07:B}
For every lattice box $\Omega$ and every $A\in\R^{\Omega^*}$,
\[
\|d_1A\|^2_{\Omega^{**}}+\|d_0^TA\|^2=K^{\rm int}_\Omega(A)+N_\partial(A).
\]
Consequently $Q^{(3.38)}_\Omega=\frac12\big[K^{\rm int}_\Omega+N_\partial+\mu^2\|\cdot\|^2_{\Omega^*}\big]$: the modified photon form is the internal Neumann--Feynman form plus the positive boundary term $N_\partial$. (By the proof of \cref{lem:s07:R}(a), the marginal of the Neumann--Dirichlet Gaussian of $\Omega$ on the internal bonds has the precision $M^{(3.38)}-N^{\rm low}_\partial-P_{12}P_{22}^{-1}P_{21}$: the modified form minus the squares of the near-normal bonds and minus the Schur term of the far-face variables.)
\end{lemma}

\begin{proof}
Extend $A$ by $0$ to a $1$-form $a$ on $\Z^d$. For finitely supported $1$-forms, $\sum_p|da(p)|^2+\sum_x|d^*a(x)|^2=\sum_\mu\sum_\nu\sum_x|a_\mu(x+e_\nu)-a_\mu(x)|^2$ (the Hodge Laplacian on $\Z^d$ is the componentwise Laplacian; summation by parts). On the left, $d^*a(x)$ is the masked divergence for $x\in\Omega$ and $0$ outside, since no internal bond has an end outside $\Omega$; a plaquette $p\notin\Omega^{**}$ meeting $\Omega^*$ has exactly one internal bond $b$ (two parallel internal edges put all four corners in $\Omega$; two adjacent ones put three corners in $\Omega$, and a product of intervals then contains the fourth), so $|da(p)|^2=|A_b|^2$. On the right, pairs with both bonds internal give $K^{\rm int}_\Omega$; a pair $(x,\mu),(x+e_\nu,\mu)$ with exactly one internal bond contributes $|A_b|^2$; for $\nu\ne\mu$ the partner is non-internal iff $x+e_\nu\notin\Omega$ (as $x,x+e_\mu\in\Omega$, and in a box $x+e_\nu\in\Omega\iff x+e_\mu+e_\nu\in\Omega$), \ie iff the plaquette spanned by $b$ and $e_\nu$ lies outside $\Omega^{**}$; these pairs are in bijection with the plaquettes counted on the left and cancel; for $\nu=\mu$ they are the terms of $N_\partial$.
\end{proof}

\begin{lemma}[the same box: the second inequality of King's (3.42) is reversed]\label{lem:s07:R}
Let $\Omega$ be a lattice box, $\langle\cdot\rangle^{\ND}_\Omega$ the Neumann--Dirichlet state on $X=\Omega$ (variables $A_\mu(x)$, $x\in\Omega$, all outgoing bonds; photon form $Q^{\ND}_\Omega=\frac12[K_\Omega+\mu^2\|\cdot\|^2]$ of \cref{def:s02:acube}) and $\langle\cdot\rangle^{(3.38)}_\Omega$ the modified box state of \cref{def:s07:box} (variables $\Omega^*$); both carry the same Dirichlet matter on $\Omega$ (constant $\dm_\eps$, free Wick constant), coupled to $A$ through the internal bonds only. Let $\Sigma^{\ND}_\Omega$ and $\Sigma^u_\Omega$ be the covariances of $F=(d_1A(p))_{p\in\Omega^{**}}$ under the two photon Gaussians. Then:
\begin{enumerate}[label=\textup{(\alph*)},nosep]
\item $\Sigma^u_\Omega=\Sigma^{(3.38)}_\Omega\le\Sigma^{\ND}_\Omega$ on $\R^{\Omega^{**}}$;
\item for every gauge-invariant $Y\in\cI$ with sources in $\Omega$ and every $e\ge0$: $\langle Y\rangle^{\ND}_\Omega\le\langle Y\rangle^{(3.38)}_\Omega$, and $\langle X\rangle^{\ND}_\Omega\ge\langle X\rangle^{(3.38)}_\Omega$ for $X\in\cD$. This is the reverse of the second inequality of \cite[(3.42)]{K2}, whose middle term is the modified box state on $C_{J+1}$; it holds weakly, for the modified state with the constant $\dm_\eps$ (nothing is asserted for King's printed site-dependent counterterm);
\item King's inequality fails strictly in the Gaussian sector: at $e=0$ and $Y=e^{\mathrm it\langle v,F\rangle}$ it reads $v^T\Sigma^{\ND}_\Omega v\le v^T\Sigma^u_\Omega v$, while $\Delta_\Omega\defeq\Sigma^{\ND}_\Omega-\Sigma^u_\Omega\ge0$ is nonzero. In exact rational arithmetic the inertia (negative, zero, positive) of $\Delta_\Omega$ is $(0,1,8)$ on the $4\times4$ box ($d=2$) for $\mu^2\in\{\frac1{20},\frac14,1\}$ and $(0,3,13)$ on the $4\times2\times2$ box ($d=3$) for $\mu^2\in\{\frac1{20},\frac14\}$; single-plaquette certificates are $\Delta_{pp}=\Var^{\ND}F(p)-\Var^{(3.38)}F(p)=\frac{14285248}{44314767}$ at $p=((0,0);0,1)$ on $4\times4$, $\mu^2=\frac14$, and $\frac{121897601600}{395580222961}$ at $p=((0,0,0);1,2)$ on $4\times2\times2$, $\mu^2=\frac1{20}$.
\end{enumerate}
\end{lemma}

\begin{proof}
(a) Write the Neumann--Dirichlet precision $P\defeq K_\Omega+\mu^2$ in blocks, the internal variables $\Omega^*$ first and the outgoing boundary variables second. Its internal block is $P_{11}=K^{\rm int}_\Omega+N^{\rm up}_\partial+\mu^2$: pairs of $K_\Omega$ with both bonds internal give $K^{\rm int}_\Omega$; a pair $(x,\mu),(x+e_\nu,\mu)$ with exactly one internal bond is impossible for $\nu\ne\mu$ (box convexity, as in \cref{lem:s07:B}) and for $\nu=\mu$ is $(x,\mu)$ internal with $(x+e_\mu,\mu)$ outgoing, \ie $x+2e_\mu\notin\Omega$ (the other assignment would need $x+e_\mu\notin\Omega$ and $x+e_\mu\in\Omega$ at once); its diagonal contribution to the internal block is $|A_\mu(x)|^2$, which is $N^{\rm up}_\partial$. Hence, by \cref{lem:s07:B}, $M^{(3.38)}\defeq K^{\rm int}_\Omega+N_\partial+\mu^2=P_{11}+N^{\rm low}_\partial\ge P_{11}$. The marginal covariance of the internal variables under the Neumann--Dirichlet Gaussian is the inverse of the Schur complement, $(P^{-1})_{11}=(P_{11}-P_{12}P_{22}^{-1}P_{21})^{-1}\ge P_{11}^{-1}\ge(M^{(3.38)})^{-1}$, by operator monotonicity of the inverse on positive definite matrices and $P_{12}P_{22}^{-1}P_{21}\ge0$. Since $F=d_1A$ on $\Omega^{**}$ uses internal bonds only, $\Sigma^{\ND}_\Omega=D_1(P^{-1})_{11}D_1^T\ge D_1(M^{(3.38)})^{-1}D_1^T=\Sigma^{(3.38)}_\Omega$, and $\Sigma^{(3.38)}_\Omega=\Sigma^u_\Omega$ by Step 0 in the proof of \cref{lem:s07:U} (the field strength depends only on the $\ker d_0^T$-component, whose law is the same under both Gaussians, Step 1 of \cref{lem:s07:RPbox}). (b) The mechanism of the second part of \cref{lem:s07:U} with $B=\Omega$ and no hopping removal: both states have the Dirichlet matter of $\Omega$ coupled to $\Omega^*$, and for gauge-invariant $Y$ with sources in $\Omega$ the expectation depends on the photon law only through the centred Gaussian law of $F$ on $\operatorname{ran}d_1$; the two $F$-laws are ordered by (a); realise them, regularise, apply \cref{lem:s02:BFS}(a) (resp.\ (d)) by interpolation and let the regularisation go to zero exactly as there. The larger covariance has the smaller $\cI$-expectation, so $\langle Y\rangle^{\ND}_\Omega\le\langle Y\rangle^{(3.38)}_\Omega$. (c) At $e=0$ the matter decouples from $A$ and is identical in the two states, so for $Y=e^{\mathrm it\langle v,F\rangle}$ both sides equal $\exp(-\frac{t^2}2v^T\Sigma v)$ with the respective $\Sigma$; King's $\le$ becomes $v^T\Sigma^{\ND}_\Omega v\le v^T\Sigma^u_\Omega v$, which is false for every $v$ with $v^T\Delta_\Omega v>0$, for instance $v=e_p$ at the plaquettes listed. The inertias and the certificates are obtained from the two covariance matrices in exact rational arithmetic (symmetric elimination with diagonal pivoting over $\mathbb Q$).
\end{proof}

\emph{Scope.} (i) King's passage from the box states to the cube limit on \cite[p.~336]{K2} needs the sandwich for the box functionals with purely imaginary curvature coefficients, \ie for the whole class $\cI$; (c) refutes the second inequality on that class exactly. (ii) For the scalar observable $\Theta(0,2h_+)$ of the printed display, (b) shows that the inequality can hold only with equality; at $e=0$ it does (identical decoupled matter), and at $e>0$ a strict failure is to be expected (the Wick square couples to $F$ at order $e^2$) but is not certified here. Both clauses are for the modified state with the constant $\dm_\eps$: if King's site-dependent counterterm differs from the free-boundary number, the two matter measures differ, the mechanism of (b) does not apply, and neither clause is proved for his printed objects; (c) alone is independent of the counterterm, the matter decoupling from $A$ at $e=0$.

\begin{remark}[what is false in the printed route]\label{rem:s07:R1}
(i) King's Neumann--Dirichlet cube state is not reflection symmetric about any plane: the site set $C_J\cap\Lk$ is symmetric in direction $\nu$ about $x_\nu=-\eps_k/2$, the set of $\nu$-bonds about $x_\nu=0$, and no reflection maps the variable set to itself (\cref{lem:s02:refl} accordingly obtains the reflection symmetry of the limit from a cut sandwich, not from a lattice symmetry). (ii) The natural symmetrisation, the Feynman gradient form on the internal bonds $K^{\rm int}_\Omega+\mu^2$ (the modified form without $N_\partial$; reflection symmetric), is \emph{not} reflection positive on gauge-invariant functionals. This is exact at $e=0$: on the $4\times4$ box at $\mu^2=\frac1{20}$ the reflection kernel $R_{pq}=\sigma_p\Sigma[\rho p,q]$ of the plus-side plaquettes is negative definite (inertia $(3,0,0)$), and on $4\times2\times2$ at $\mu^2=\frac1{20}$ it has inertia $(2,2,2)$ with $R_{pp}=\Cov(\theta F(p),F(p))<0$ for $p=((2,0,0);0,1)$. Indeed, at $e=0$ the matter decouples, and the gauge-invariant plus-side functionals $G_1=e^{\mathrm it\langle v,F_+\rangle}$, $G_2=1$ give the matrix $(\langle\theta(G_i)G_j\rangle)$ with determinant $e^{-t^2v^T\Sigma_{++}v}(e^{t^2v^TRv}-1)<0$ for every $t>0$ when $v^TRv<0$, since $\langle\theta(G_1)G_1\rangle=\exp(-t^2v^T\Sigma_{++}v+t^2v^TRv)$ and $\langle\theta(G_1)\rangle=\langle G_1\rangle=\exp(-\frac12t^2v^T\Sigma_{++}v)$ by the reflection symmetry $\Sigma_{--}=\Sigma_{++}$. By continuity at fixed spacing the symmetrised cube state is not reflection positive for small $e>0$ either. A proof must work at every coupling, so reflection positivity of King's own cube form, or of this symmetrisation, at fixed spacing is not available as a route: any route through cube states at fixed spacing has to change the photon form, as the modified box does (a reflection-symmetric, reflection-positive modification whose values converge to the same limit, \cref{thm:s08:OS2C}). Nothing is claimed about the symmetrised form at a given larger $e$ or about other proofs; the route taken here needs the comparison of \cref{lem:s07:L,lem:s07:U}, and its reflecting plane must be a lattice mid-plane for the half-bond mechanism of \cref{lem:s07:RPbox} (King's cubes have an even number of layers). (iii) Accordingly, the second inequality of \cite[(3.42)]{K2} is false as printed in the following sense: false for the class of observables that the argument on \cite[p.~336]{K2} needs (exactly, \cref{lem:s07:R}(c)), and never strictly true for the modified state with the constant counterterm (\cref{lem:s07:R}(b)). Its valid first half is \cref{lem:s07:L}; the invalid half is replaced by \cref{lem:s07:U} with the margin $R(\lmu)$ of \eqref{eq:s07:margin}, $n(\lmu,c)$ cube steps instead of one.
\end{remark}

\begin{remark}[the box form is the larger one]\label{rem:s07:K8}
King describes the modified box action as obtained from the cube action by omitting some terms. By \cref{lem:s07:B} the relation is the reverse: on one box the modified photon form is the internal Neumann--Dirichlet form $P_{11}$ plus the positive term $N^{\rm low}_\partial$ and without the outgoing boundary variables, so the modified state is the one with the \emph{larger} form, the smaller field-strength covariance, and the larger $\cI$-expectations: $\langle Y\rangle^{\ND}_{C}\le\langle Y\rangle^{(3.38)}_{C}$ on the cone for every box $C$ (\cref{lem:s07:R}(b)), the direction that \cref{lem:s02:cut} requires for consistency with \cref{lem:s07:L,lem:s07:U}. For King's convention the same comparison holds with Step~1 of \cref{lem:s04:G} prepended (\cref{rem:s04:F}): $\avK{Y}\le\ava{Y}\le\langle Y\rangle^{(3.38)}_{k,\XJ}$ on $\cI$.
\end{remark}

\subsection{The sandwich}

\begin{lemma}[the lower half]\label{lem:s07:L}
Let $C\subset\Omega$ be lattice boxes with $C$ plus one layer contained in $\Omega$, so that the variables $A_\mu(x)$, $x\in C$, of the Neumann--Dirichlet state on $C$ (the outgoing bonds) are internal bonds of $\Omega$. Then
\begin{enumerate}[label=\textup{(\alph*)},nosep]
\item $Q^{(3.38)}_\Omega\ge Q^{\ND}_C\oplus\frac{\mu^2}2I_{\rm rest}$ as quadratic forms on $\R^{\Omega^*}=\R^{C\text{-variables}}\oplus\R^{\rm rest}$;
\item for sources such that every plaquette and every site at which they do not vanish lies in $C$,
\[
\langle Y\rangle^{\ND}_C\le\langle Y\rangle^{(3.38)}_\Omega\ \ (Y\in\cI),\qquad\langle X\rangle^{\ND}_C\ge\langle X\rangle^{(3.38)}_\Omega\ \ (X\in\cD),
\]
and the same for $Z^a$ (equal positive multipliers, \cref{lem:s02:alg}(a)). The $\cD$ half uses \cref{lem:s02:BFS}(d); the $\cI$ half only (a), (b), (e).
\end{enumerate}
\end{lemma}

\begin{proof}
(a) $K^{\rm int}_\Omega$ contains every gradient pair of $K_C$ (both base points in $C$, both bonds internal to $\Omega$); the remaining pairs and $N_\partial$ are nonnegative; the mass term splits. (b) Start from $\langle\cdot\rangle^{(3.38)}_\Omega$ and remove the hopping cosines across $\partial C$ and all cosines not internal to $C$; by \cref{lem:s02:BFS}(b),(d) $\langle Y\rangle$ does not increase and $\langle X\rangle$ does not decrease, every intermediate measure being a BFS measure. The matter is now the Dirichlet matter of $C$: the diagonal parts $\frac12r^2$ of the cut bonds stay and are exactly the Dirichlet terms, as in the proof of \cref{lem:s02:cut}; it is coupled to $A$ only through bonds internal to $C$, which are $C$-variables, and the matter of $\Omega\setminus C$ is decoupled and cancels (\cref{lem:s02:BFS}(e)). Then pass from the form $Q^{(3.38)}_\Omega$ to the smaller form $Q^{\ND}_C\oplus\frac{\mu^2}2I_{\rm rest}$: by (a) the difference is positive semidefinite, a finite sum $\sum_is_i|\ell_i(A)|^2$ with $s_i\ge0$, and adding each term does not decrease $\langle Y\rangle$ and does not increase $\langle X\rangle$ (\cref{lem:s02:BFS}(a),(d)). So the state with the larger form, the modified one, has the larger $\langle Y\rangle$ and the smaller $\langle X\rangle$. The state with the smaller form factorises as $\langle\cdot\rangle^{\ND}_C$ times an independent Gaussian on the rest, which cancels.
\end{proof}

For King's convention the lower end is used in the form $\avK{Y}\le\ava{Y}\le\langle Y\rangle^{(3.38)}_{\Omega}$ on $\cI$ (Step~1 of \cref{lem:s04:G}, then \cref{lem:s07:L}; \cref{rem:s04:F}), and the upper end needs no change (\cref{lem:s04:D}).

\begin{lemma}[the upper half, with a margin]\label{lem:s07:U}
Let $\Omega\subset B$ be lattice boxes at spacing $\eps$ such that $B$ contains every site within $\ell^\infty$-distance $N$ lattice steps of $\Omega$, and let $\mu=\lmu\eps$ satisfy $d(\pi/2N)^2\le\mu^2$, \ie
\[
N\eps\ \ge\ \frac{\pi\sqrt d}{2\lmu}\qquad\big(=R(\lmu)=\tfrac{\pi\sqrt3}{2\lmu}\text{ in }d=3\big).
\]
Let $\Sigma_X$ denote the covariance of $(\eps F(p))_{p\in\Omega^{**}}$ under the photon Gaussian of the state on $X$. Then
\begin{enumerate}[label=\textup{(\alph*)},nosep]
\item $\Sigma^{\ND}_B\le\Sigma^{(3.38)}_\Omega=\Sigma^u_\Omega$ on $\R^{\Omega^{**}}$;
\item for sources supported in $\Omega$ as in \cref{lem:s07:L},
\[
\langle Y\rangle^{(3.38)}_\Omega\le\langle Y\rangle^{\ND}_B\ \ (Y\in\cI),\qquad\langle X\rangle^{(3.38)}_\Omega\ge\langle X\rangle^{\ND}_B\ \ (X\in\cD),
\]
and the same for $Z^a$. Only the $\cI$ half is used in the sequel.
\end{enumerate}
\end{lemma}

\begin{proof}
\emph{(a), in lattice units.} Step 0: $\Sigma^{(3.38)}_\Omega=\Sigma^u_\Omega$, because $d_1A=d_1A^T$ depends only on the $\ker d_0^T$-component, whose law $\propto e^{-Q^u_\Omega(A^T)}$ is the same under both Gaussians (Step 1 of \cref{lem:s07:RPbox}). Step 1 (variational identity): for a Gaussian with density $\propto e^{-\frac12A^TM_XA}$, $\langle g,\Sigma_Xg\rangle=\Var\langle g,d_1A\rangle=\sup_A[2\langle g,d_1A\rangle-A^TM_XA]$, the supremum attained at $A=M_X^{-1}d_1^Tg$. So it suffices to show that for every $A$ on the $B$-variables there is $A'\in\R^{\Omega^*}$ with $d_1A'=d_1A$ on $\Omega^{**}$ and $\|d_1A'\|^2_{\Omega^{**}}+\mu^2\|A'\|^2_{\Omega^*}\le K_B(A)+\mu^2\|A\|^2_B$; take $A'\defeq A|_{\Omega^*}$ (the plaquettes of $\Omega^{**}$ use only $\Omega^*$-bonds). Step 2 (cutoff): let $\psi(x)=\prod_\nu\phi_\nu(x_\nu)$ with $\phi_\nu=1$ on the $\nu$-range of $\Omega$ and on the first layer outside, $\phi_\nu(t)=\cos(\pi(t-1)/2N)$ at the $t$-th layer outside for $1\le t\le N+1$ (so $\phi_\nu(N+1)=0$), and $0$ beyond; then $0\le\psi\le1$, $\psi=1$ on $\Omega$ and its neighbours, $\supp\psi\subset B$. Put $a\defeq\psi A$ on $\Z^d$ (the bond $\langle x,x+e_\mu\rangle$ weighted by $\psi(x)$; $A\defeq0$ off the $B$-variables). Step 3 (Hodge on $\Z^d$): $\|d_1A\|^2_{\Omega^{**}}=\|d_1A'\|^2_{\Omega^{**}}\le\sum_p|da(p)|^2\le\sum_p|da|^2+\sum_x|d^*a|^2=\sum_\mu\sum_\nu\sum_x|a_\mu(x+e_\nu)-a_\mu(x)|^2$, since the bonds of plaquettes in $\Omega^{**}$ have base points in $\Omega$, where $\psi=1$. Step 4 (discrete IMS identity): for finitely supported $\psi f$ and a direction $e$,
\[
\sum_x|\psi(x+e)f(x+e)-\psi(x)f(x)|^2=\sum_x\psi(x)\psi(x+e)|f(x+e)-f(x)|^2+\sum_x\psi(x)\big(2\psi(x)-\psi(x+e)-\psi(x-e)\big)f(x)^2 ,
\]
exactly (per bond the difference of the two sides is $\psi'(\psi'-\psi)f'^2+\psi(\psi-\psi')f^2$; collect by site). With $f=A_\mu$ and $e=e_\nu$: $\psi\psi'\le1$, and $\psi\psi'>0$ only for pairs with both base points in $B$, which are gradient pairs of $K_B$; so the first sums, over $\mu,\nu$, total at most $K_B(A)$. The potential is $V(x)=\sum_\nu\big(\prod_{\kappa\ne\nu}\phi_\kappa^2\big)\phi_\nu(-\Delta\phi_\nu)(x_\nu)$. Step 5 (the profile): with $\vartheta\defeq\pi/2N$, $\phi(-\Delta\phi)=0$ on the range of $\Omega$ (all neighbours equal $1$), $=1-\cos\vartheta\le\vartheta^2/2$ at $t=1$, $=2\cos^2(\vartheta(t-1))(1-\cos\vartheta)\le\vartheta^2$ for $2\le t\le N$ (using $\phi(N+1)=\cos(\vartheta N)=0$ at $t=N$), and $=0$ for $t\ge N+1$; hence $V=0$ on $\Omega$ and $V\le d\vartheta^2$ everywhere. Step 6: $\|d_1A'\|^2_{\Omega^{**}}\le K_B(A)+d\vartheta^2\sum_{x\in B\setminus\Omega}|A(x)|^2$; the variables with base point in $B\setminus\Omega$ and those of $\Omega^*$ (base point in $\Omega$) are disjoint subsets of the $B$-variables, so $\|d_1A'\|^2_{\Omega^{**}}+\mu^2\|A'\|^2_{\Omega^*}\le K_B(A)+\mu^2\|A\|^2_B$ as soon as $d\vartheta^2\le\mu^2$.

\emph{(b).} Start from $\langle\cdot\rangle^{\ND}_B$ and remove all hopping cosines touching $B\setminus\Omega$ (\cref{lem:s02:BFS}(b): $\cI$ does not increase; (d): $\cD$ does not decrease). The matter is now the Dirichlet matter of $\Omega$ (the diagonals of the bonds leaving $\Omega$ kept; every such bond ends in $B$ since $N\ge1$), coupled to $A$ on $\Omega^*$, and the sites of $B\setminus\Omega$ decouple. This matter is identical to that of the modified box state: the same cosines on $\Omega^*$, the same diagonals $\frac12r_x^2$ for every bond leaving $\Omega$, the same single-site weights with the constant $\dm_\eps$ and the same Wick constant. For a gauge-invariant observable with sources in $\Omega$ the matter integral at fixed $A$, $\Psi_Y(A)=\int\prod_{x\in\Omega}d^2\phi_x\,Y\,e^{-S_m}$, is invariant under $A\mapsto A+d_0\chi$, $\chi\in\R^\Omega$ (as in Step 1 of \cref{lem:s07:RPbox}), hence a function of the class of $A|_{\Omega^*}$ modulo $\operatorname{ran}d_0=\ker d_1$ (a box has $H^1=0$: the cycle space of a grid graph is generated by its $4$-cycles), hence a function $\tilde\Psi_Y(d_1A|_{\Omega^{**}})$ of the field strengths on $\Omega^{**}$. So $\langle Y\rangle=\E[\tilde\Psi_Y(F)]/\E[\tilde\Psi_1(F)]$ depends on the photon law only through the centred Gaussian law of $F=(d_1A(p))_{p\in\Omega^{**}}$, with covariance $\Sigma^{\ND}_B$ on the Neumann--Dirichlet side after the removal and $\Sigma^u_\Omega$ on the modified side (Step 0). Both laws live on $\operatorname{ran}d_1$. Let $T:\R^{\Omega^{**}}\to\R^{\Omega^*}$ be a linear map with $d_1T=\mathrm{id}$ on $\operatorname{ran}d_1$, extended linearly to all of $\R^{\Omega^{**}}$ (any extension serves, both laws being supported on $\operatorname{ran}d_1$), and realise each law as the photon Gaussian of $A''\defeq TF\in\R^{\Omega^*}$, covariance $T\Sigma T^T$: then $d_1A''=F$ almost surely, the gauge-invariant expectations are reproduced, and $T\Sigma^{\ND}_BT^T\le T\Sigma^u_\Omega T^T$ by (a). These covariances are degenerate; regularise to $T\Sigma T^T+\delta I$, $\delta>0$, which are non-degenerate, so the regularised states are BFS measures, the ordering of the covariances is preserved, and the law of $F''=d_1A''$ is $N(0,\Sigma+\delta d_1d_1^T)$, still ordered. \Cref{lem:s02:BFS}(a) by interpolation (the difference of the precision matrices is positive semidefinite by operator monotonicity of the inverse, hence a finite sum of terms $s|\ell(A)|^2$, $s\ge0$) gives $\langle Y\rangle_{\Sigma^u+\delta}\le\langle Y\rangle_{\Sigma^{\ND}+\delta}$ for $Y\in\cI$, and the reverse for $\cD$ by \cref{lem:s02:BFS}(d). \emph{The limit $\delta\downarrow0$.} For the $\cI$ half, $\tilde\Psi_Y$ is bounded and continuous: $|\Psi_Y(A)|\le\int e^{r^2(g)}e^{\sum\eps^{d-2}r_xr_y}\prod_xd\rho_x<\infty$ independently of $A$ ($|e^{\mathrm iF}|=1$; the hopping bound $e^{\eps^{d-2}r_xr_y}\le e^{\frac12\eps^{d-2}(r_x^2+r_y^2)}$ cancels the diagonals and $\lambda>0$ does the rest), and so is $\tilde\Psi_1$, with $\tilde\Psi_1>0$; the Gaussians $N(0,\Sigma+\delta d_1d_1^T)$ converge weakly to $N(0,\Sigma)$, so numerator and denominator converge and the inequality passes to the limit. For the $\cD$ half (not used below), $X=\Theta(f,-g)$ gives $\tilde\Psi_X(F)=e^{\langle v,F\rangle}\tilde\Psi^0_X(F)$ with $v$ the plaquette vector of $f$ and $0<\tilde\Psi^0_X\le e^{2\sW_\eps\sS(g)}\tilde\Psi_1\le C$ bounded and continuous, so $\tilde\Psi_X$ is continuous but unbounded; under $N(0,\Sigma+\delta d_1d_1^T)$ one has $\E[e^{2\langle v,F''\rangle}]=e^{2v^T(\Sigma+\delta d_1d_1^T)v}\le e^{2v^T\Sigma v+2\|d_1^Tv\|^2}$ for $\delta\le1$, so $\{\tilde\Psi_X(F'')\}_{0<\delta\le1}$ is bounded in $L^2$, hence uniformly integrable, and weak convergence with uniform integrability and continuity gives $\E[\tilde\Psi_X(F'')]\to\E[\tilde\Psi_X(F)]$. Altogether $\langle Y\rangle^{(3.38)}_\Omega\le\langle Y\rangle^{\text{after removal}}\le\langle Y\rangle^{\ND}_B$, and the reverse for $\cD$.
\end{proof}

\begin{remark}[why a margin, and why a fixed photon mass]
Part (a) of \cref{lem:s07:U} fails without a margin: for $B=\Omega$ the comparison is reversed (\cref{lem:s07:R}), and with a single layer ($N=1$) the difference $\Sigma^u_\Omega-\Sigma^{\ND}_B$ has a negative eigenvalue already at moderate $\mu$. No fixed margin works at $\lmu=0$: there $\Sigma^u_\Omega$ is the projector onto exact $2$-forms, while $\Var_B\langle g,dA\rangle>\|g\|^2$ for a Coulomb potential whose gradient leaks out of $B$. So \cref{lem:s07:U} is a statement at fixed $\lmu>0$ with $R(\lmu)\to\infty$ as $\lmu\to0$; a limit of vanishing photon mass inherits reflection positivity, if it exists, as a pointwise limit of positive semidefinite matrices. On one box the three states are ordered on $\cI$ as $\langle Y\rangle^{\ND}_\Omega\le\langle Y\rangle^{(3.38)}_\Omega\le\langle Y\rangle^{\ND}_B$ (\cref{lem:s07:R,lem:s07:U}), as the monotonicity of \cref{lem:s02:cut} requires. The sandwich with the boxes \eqref{eq:s07:margin} is assembled in the proof of \cref{thm:s08:OS2C}; for King's convention its lower end carries the kill step of \cref{rem:s04:F} and its upper end is unchanged (\cref{lem:s04:D}).
\end{remark}
